\section{Semicircles, constant drivers, and a restricted rigidity theorem}\label{sec:geometry}

Choose a plane containing a direction tangent to the defect and its
normal, and write $z=x+\ii y$ with $y>0$. The common-reference semicircle
from $r>0$ to zero is
\begin{equation}\label{eq:semicircle}
 z_r(\theta)=\frac r2(1+\cos\theta+\ii\sin\theta),\qquad0\leq\theta<\pi.
\end{equation}
Inversion preserves the upper half-plane and gives
\begin{equation}
 w_r(\theta)=-\frac1{z_r(\theta)}
 =-\frac1r+\frac{\ii}{r}\tan(\theta/2).
\end{equation}
This is a vertical slit. In the chordal convention
\begin{equation}\label{eq:loewner}
 \partial_t g_t(w)=\frac2{g_t(w)-U_t},\qquad g_0(w)=w,
\end{equation}
the constant driver $U_t=a=-1/r$ has solution
\begin{equation}\label{eq:constant-driver}
 g_t(w)=a+\sqrt{(w-a)^2+4t},\quad
 \gamma_t=a+2\ii\sqrt t,\quad z_t=-\frac1{a+2\ii\sqrt t}.
\end{equation}
The square-root branch is fixed by $g_t(w)\sim w$ at infinity. Its
expansion is $g_t(w)=w+2t/w+O(w^{-2})$, so the half-plane capacity is
$2t$, consistent with \cite[Lecture~2, Proposition~2.17]{Lawler}.
Comparing the traces gives the exact dictionary
\begin{equation}\label{eq:capacity-dictionary}
 t=\frac{\tan^2(\theta/2)}{4r^2}.
\end{equation}
The complete semicircle reaches the reference only as $t\to\infty$.
At fixed angle, $r\mapsto\lambda r$ sends $t\mapsto\lambda^{-2}t$.
Thus this inverted capacity is neither the coordinate $\log r$ nor
the arithmetic shift $\omega$. The formulas identify contours; they
do not assert equality of renormalized expectations through an
inversion singular at an endpoint.

\subsection{Fixed scalar coupling and fixed special charge}

Let the Euclidean Clifford generators satisfy
$\{\Gamma_A,\Gamma_B\}=2\delta_{AB}$. Take one internal direction $I$
and fix a nonzero complex spinor $\epsilon_c$ obeying
\begin{equation}\label{eq:inherited-charge}
 \epsilon_s=0,\qquad J\epsilon_c=-\epsilon_c,
 \qquad J=\ii\Gamma_I\Gamma_3.
\end{equation}
Here $\epsilon_s+X^\mu\Gamma_\mu\epsilon_c$ is the relevant conformal
Killing spinor, with its constant part set to zero. These are the
inherited common bulk conditions. For a regular real curve
$X=(x,y)\ne0$ in the $(1,3)$-plane and constant positive scalar
arclength coupling, its local line condition is
\begin{equation}\label{eq:rigidity-condition}
 (\ii\dot x\Gamma_1+\ii\dot y\Gamma_3+v\Gamma_I)
 (x\Gamma_1+y\Gamma_3)\epsilon_c=0,
 \qquad v=|\dot X|>0.
\end{equation}

\begin{proposition}[Rigidity under \eqref{eq:inherited-charge}]\label{prop:rigidity}
Equation~\eqref{eq:rigidity-condition} holds if and only if
\begin{equation}\label{eq:rigid-tangent}
 \frac{\dot z}{|\dot z|}=\frac{\ii z^2}{|z|^2}.
\end{equation}
Consequently $-1/z$ has constant real part and increasing imaginary
part. Connected solutions in the upper half-plane are subarcs of
circles through zero with center on the real axis, or the degenerate
imaginary-axis line. A solution from a nonzero real boundary point to
zero is the corresponding oriented upper semicircle.
\end{proposition}
\begin{proof}
The projector implies $\Gamma_I\epsilon_c=-\ii\Gamma_3\epsilon_c$.
Clifford multiplication reduces \eqref{eq:rigidity-condition} to
$(A+B\Gamma_1\Gamma_3)\epsilon_c=0$, where
\begin{equation}
 A=x\dot x+y\dot y+vy,\qquad B=y\dot x-x\dot y+vx.
\end{equation}
Since $(\Gamma_1\Gamma_3)^2=-1$, multiplication by
$A-B\Gamma_1\Gamma_3$ gives $(A^2+B^2)\epsilon_c=0$.
The coefficients are real and the spinor is nonzero, so $A=B=0$.
Solving gives
\begin{equation}
 \frac{\dot x}{v}=-\frac{2xy}{x^2+y^2},\qquad
 \frac{\dot y}{v}=\frac{x^2-y^2}{x^2+y^2},
\end{equation}
which is \eqref{eq:rigid-tangent}. Conversely it makes $A=B=0$.
Moreover $\mathrm d(-1/z)/\mathrm ds=\ii v/|z|^2$.
The level set $\Ree(-1/z)=a\ne0$ is
$x^2+y^2=-x/a$, and $a=0$ gives the imaginary axis.
\end{proof}

For instance, an inverted deformation
$w(s)=a+\varepsilon h(s)+\ii s$ preserves these conditions only if
$\varepsilon h'(s)=0$. A constant $h$ changes the slit location;
a nontrivial fixed-endpoint deformation fails. A regular
capacity-parametrized Loewner trace satisfying the same conditions
therefore has constant driving.

This is a bulk theorem for one plane, one charge and one scalar
prescription. It is not a classification of supersymmetric contours
with varying scalar directions, other conformal Killing spinors,
additional endpoint degrees of freedom, or nonplanar embeddings.
Section~\ref{sec:projectors} supplies a constructive control showing
that a changed scalar prescription allows arbitrary tangents.
